%%% The main file. It contains definitions of basic parameters and includes all other parts.

%% Settings for single-side (simplex) printing
% Margins: left 40mm, right 25mm, top and bottom 25mm
% (but beware, LaTeX adds 1in implicitly)
\documentclass[12pt,a4paper]{report}
\setlength\textwidth{145mm}
\setlength\textheight{247mm}
\setlength\oddsidemargin{15mm}
\setlength\evensidemargin{15mm}
\setlength\topmargin{0mm}
\setlength\headsep{0mm}
\setlength\headheight{0mm}
% \openright makes the following text appear on a right-hand page
\let\openright=\clearpage

%% Settings for two-sided (duplex) printing
% \documentclass[12pt,a4paper,twoside,openright]{report}
% \setlength\textwidth{145mm}
% \setlength\textheight{247mm}
% \setlength\oddsidemargin{14.2mm}
% \setlength\evensidemargin{0mm}
% \setlength\topmargin{0mm}
% \setlength\headsep{0mm}
% \setlength\headheight{0mm}
% \let\openright=\cleardoublepage

%% Generate PDF/A-2u
\usepackage[a-2u]{pdfx}

%% Character encoding: usually latin2, cp1250 or utf8:
\usepackage[utf8]{inputenc}

%% Prefer Latin Modern fonts
\usepackage{lmodern}

%% Further useful packages (included in most LaTeX distributions)
\usepackage{amsmath}        % extensions for typesetting of math
\usepackage{amsfonts}       % math fonts
\usepackage{amsthm}         % theorems, definitions, etc.
\usepackage{bbding}         % various symbols (squares, asterisks, scissors, ...)
\usepackage{bm}             % boldface symbols (\bm)
\usepackage{graphicx}       % embedding of pictures
\usepackage{fancyvrb}       % improved verbatim environment
\usepackage{natbib}  
\setcitestyle{authordate}
      % citation style AUTHOR (YEAR), or AUTHOR [NUMBER]
\usepackage[nottoc]{tocbibind} % makes sure that bibliography and the lists
			    % of figures/tables are included in the table
			    % of contents
\usepackage{dcolumn}        % improved alignment of table columns
\usepackage{booktabs}       % improved horizontal lines in tables
\usepackage{paralist}       % improved enumerate and itemize
%\usepackage[usenames]{xcolor}  % typesetting in color
%%% Basic information on the thesis

% Thesis title in English (exactly as in the formal assignment)
\def\ThesisTitle{Czech NLP Processing with Contextualized Embeddings}

% Author of the thesis
\def\ThesisAuthor{Bc. Petra Vysušilová}

% Year when the thesis is submitted
\def\YearSubmitted{2020}

% Name of the department or institute, where the work was officially assigned
% (according to the Organizational Structure of MFF UK in English,
% or a full name of a department outside MFF)
\def\Department{Institute of Formal and Applied Linguistics}

% Is it a department (katedra), or an institute (ústav)?
\def\DeptType{Institute}

% Thesis supervisor: name, surname and titles
\def\Supervisor{RNDr. Milan Straka, Ph.D.}

% Supervisor's department (again according to Organizational structure of MFF)
\def\SupervisorsDepartment{Institute of Formal and Applied Linguistics}

% Study programme and specialization
\def\StudyProgramme{Computer science}
\def\StudyBranch{Artificial intelligence}

% An optional dedication: you can thank whomever you wish (your supervisor,
% consultant, a person who lent the software, etc.)
\def\Dedication{%
Dedication.
}

% Abstract (recommended length around 80-200 words; this is not a copy of your thesis assignment!)
\def\Abstract{%
Recently, several methods for unsupervised pre-training of contextualized word embeddings have been proposed, most importantly the BERT model (Devlin et al., 2018). Such contextualized representations have been extremely useful as additional features in many NLP tasks like morphosyntactic analysis, entity recognition or text classification.

Most of the evaluation have been carried out on English. However, several of the released models have been pre-trained on many languages including Czech, like multilingual BERT or XLM-RoBERTa (Conneau et al, 2019). Therefore, the goal of this thesis is to perform experiments quantifying improvements of employing pre-trained contextualized representation in Czech natural language processing.
}

% 3 to 5 keywords (recommended), each enclosed in curly braces
\def\Keywords{%
{key} {words}
}

%% The hyperref package for clickable links in PDF and also for storing
%% metadata to PDF (including the table of contents).
%% Most settings are pre-set by the pdfx package.
\hypersetup{unicode}
\hypersetup{breaklinks=true}

% Definitions of macros (see description inside)
%%% This file contains definitions of various useful macros and environments %%%
%%% Please add more macros here instead of cluttering other files with them. %%%

%%% Minor tweaks of style

% These macros employ a little dirty trick to convince LaTeX to typeset
% chapter headings sanely, without lots of empty space above them.
% Feel free to ignore.
\makeatletter
\def\@makechapterhead#1{
  {\parindent \z@ \raggedright \normalfont
   \Huge\bfseries \thechapter. #1
   \par\nobreak
   \vskip 20\p@
}}
\def\@makeschapterhead#1{
  {\parindent \z@ \raggedright \normalfont
   \Huge\bfseries #1
   \par\nobreak
   \vskip 20\p@
}}
\makeatother

% This macro defines a chapter, which is not numbered, but is included
% in the table of contents.
\def\chapwithtoc#1{
\chapter*{#1}
\addcontentsline{toc}{chapter}{#1}
}

% Draw black "slugs" whenever a line overflows, so that we can spot it easily.
\overfullrule=1mm

%%% Macros for definitions, theorems, claims, examples, ... (requires amsthm package)

\theoremstyle{plain}
\newtheorem{thm}{Theorem}
\newtheorem{lemma}[thm]{Lemma}
\newtheorem{claim}[thm]{Claim}

\theoremstyle{plain}
\newtheorem{defn}{Definition}

\theoremstyle{remark}
\newtheorem*{cor}{Corollary}
\newtheorem*{rem}{Remark}
\newtheorem*{example}{Example}

%%% An environment for proofs

%%% FIXME %%% \newenvironment{proof}{
%%% FIXME %%%   \par\medskip\noindent
%%% FIXME %%%   \textit{Proof}.
%%% FIXME %%% }{
%%% FIXME %%% \newline
%%% FIXME %%% \rightline{$\square$}  % or \SquareCastShadowBottomRight from bbding package
%%% FIXME %%% }

%%% An environment for typesetting of program code and input/output
%%% of programs. (Requires the fancyvrb package -- fancy verbatim.)

\DefineVerbatimEnvironment{code}{Verbatim}{fontsize=\small, frame=single}

%%% The field of all real and natural numbers
\newcommand{\R}{\mathbb{R}}
\newcommand{\N}{\mathbb{N}}

%%% Useful operators for statistics and probability
\DeclareMathOperator{\pr}{\textsf{P}}
\DeclareMathOperator{\E}{\textsf{E}\,}
\DeclareMathOperator{\var}{\textrm{var}}
\DeclareMathOperator{\sd}{\textrm{sd}}

%%% Transposition of a vector/matrix
\newcommand{\T}[1]{#1^\top}

%%% Various math goodies
\newcommand{\goto}{\rightarrow}
\newcommand{\gotop}{\stackrel{P}{\longrightarrow}}
\newcommand{\maon}[1]{o(n^{#1})}
\newcommand{\abs}[1]{\left|{#1}\right|}
\newcommand{\dint}{\int_0^\tau\!\!\int_0^\tau}
\newcommand{\isqr}[1]{\frac{1}{\sqrt{#1}}}

%%% Various table goodies
\newcommand{\pulrad}[1]{\raisebox{1.5ex}[0pt]{#1}}
\newcommand{\mc}[1]{\multicolumn{1}{c}{#1}}


% Title page and various mandatory informational pages
\begin{document}
%%% Title page of the thesis and other mandatory pages

%%% Title page of the thesis

\pagestyle{empty}
\hypersetup{pageanchor=false}
\begin{center}

\centerline{\mbox{\includegraphics[width=166mm]{../img/logo-en.pdf}}}

\vspace{-8mm}
\vfill

{\bf\Large MASTER THESIS}

\vfill

{\LARGE\ThesisAuthor}

\vspace{15mm}

{\LARGE\bfseries\ThesisTitle}

\vfill

\Department

\vfill

\begin{tabular}{rl}

Supervisor of the master thesis: & \Supervisor \\
\noalign{\vspace{2mm}}
Study programme: & \StudyProgramme \\
\noalign{\vspace{2mm}}
Study branch: & \StudyBranch \\
\end{tabular}

\vfill

% Zde doplňte rok
Prague \YearSubmitted

\end{center}

\newpage

%%% Here should be a bound sheet included -- a signed copy of the "master
%%% thesis assignment". This assignment is NOT a part of the electronic
%%% version of the thesis. DO NOT SCAN.

%%% A page with a solemn declaration to the master thesis

\openright
\hypersetup{pageanchor=true}
\pagestyle{plain}
\pagenumbering{roman}
\vglue 0pt plus 1fill

\noindent
I declare that I carried out this master thesis independently, and only with the cited
sources, literature and other professional sources.

\medskip\noindent
I understand that my work relates to the rights and obligations under the Act No.~121/2000 Sb.,
the Copyright Act, as amended, in particular the fact that the Charles
University has the right to conclude a license agreement on the use of this
work as a school work pursuant to Section 60 subsection 1 of the Copyright Act.

\vspace{10mm}

\hbox{\hbox to 0.5\hsize{%
In ........ date ............	% FIXME!
\hss}\hbox to 0.5\hsize{%
signature of the author
\hss}}

\vspace{20mm}
\newpage

%%% Dedication

\openright

\noindent
\Dedication

\newpage

%%% Mandatory information page of the thesis

\openright

\vbox to 0.5\vsize{
\setlength\parindent{0mm}
\setlength\parskip{5mm}

Title:
\ThesisTitle

Author:
\ThesisAuthor

\DeptType:
\Department

Supervisor:
\Supervisor, \SupervisorsDepartment

Abstract:
\Abstract

Keywords:
\Keywords

\vss}

\newpage

\openright
\pagestyle{plain}
\pagenumbering{arabic}
\setcounter{page}{1}


%%% A page with automatically generated table of contents of the master thesis

\tableofcontents

%%% Each chapter is kept in a separate file
\chapter*{Introduction}
\addcontentsline{toc}{chapter}{Introduction}


This work aim to improve some natural language processing (NLP) tasks for Czech \todo{seznam zkratek} with the use of recently published artificial intelligence (AI) state-of-the-art techniques. In Deep Learning book \cite, NLP is
defined as "...the use of human languages, such as English or French, by a computer." \cite[]{Goodfellow-et-al-2016}. NLP offers a variety of problems to solve from oral--written language conversion, machine translation, syntax analysis used for automatic grammar correction as well as it serve as a base for further linguistic processing. Semantic analysis includes in addition to already mentioned machine translation tasks like sentiment analysis, natural language text generation or recognition
of homonymy or polysemy of given words. It could solve sophisticated assignments as answering questions about the input text document.

\cite{BERT_ORIG}
\section*{Tasks definition}
\todo{obrazky k taskum s priklady}
This work aims to apply the most successful NLP methods of recent years to Czech NLP task - namely tagging, lemmatization and sentiment analysis. Tagging and lemmatization represents syntax analysis, in contrast with sentiment analysis which represents semantic type of tasks. \todo{dát priklady jinych tasku}
The practical part of this work is focused on these tasks:
\begin{itemize}
\item POS tagging \\
\textit{input}: a word \\
\textit{output}: part-of-speech tags -- as noun, pronoun, punctuation mark etc.
\item lemmatization \\
\textit{input:} a word \\
\textit{output:} lemma -- a base form of a given words, meaning for example nominative of singular for nouns or infinitive for verbs. 
\item sentiment analysis \\
\textit{input:} a sentence or a sequence of sentences \\
\textit{output:} prevailing sentiment of the input from categories: neutral, positive, negative.
\todo{doplnit diskuzi o ruzných možnostech definice}
\end{itemize}.

These tasks were chosen to show how pre-trained multilingual language models can help with different types of NLP tasks in one of trained langugages. 

\section*{Text structure}
Theoretical background in NLP and used AI methods and related work is provided in the \hyperref[chap:theandme]{following chapter}.

Implementation documentation in chapter \ref{chap:impl} is followed by discussion about used methods and experiment results in chapter \ref{chap:diss} Result models are accessible to user exploration as described in chapter \ref{chap:userdoc} and text is closed by \hyperref[chap:concl]{conclusion} with future work proposals.



\chapter{Title of the first chapter}

An~example citation: \cite{Andel07}

\section{Title of the first subchapter of the first chapter}

\section{Title of the second subchapter of the first chapter}

\chapter{Implementation analysis}
\label{chap:impl}
The main purpose of this chapter is to offer the technical description of the code accompanying this work for better reproducibility and possible further experiments on every of presented tasks. This chapter describes an implementation of all language models, other related code and also presents all used libraries and technologies.
% and this chapter is concluded with a presentation of experiment types common for training of both tasks presented in following chapters.
\par
All code forms an attachment of this work and is also publicly available in GitHub \footnote{https://github.com/flower-go/DiplomaThesis}. Experiment were performed on Artificial Intelligence Cluster (AIC) \footnote{https://aic.ufal.mff.cuni.cz/} provided by Institute of Formal and Applied Linguistics, Charles University \footnote{https://ufal.mff.cuni.cz/home-page}.

\section{Code description}
This section will describe the code -- technologies and hardware used for experiments, where to find scripts for replicating experiments and how to run them.
\subsection{Technologies description}
All code is implemented in Python (v3.6.9). Python is a popular language for machine learning, because of easy use and many available libraries, which allows to focus on high-level problem solving instead of technical details. All dependencies and used libraries are listed in the \texttt{/code/requirements.txt} file, but the most important libraries should be mentioned explicitly.
\subsubsection{Tensorflow and Keras}
The main library used for developing deep learning models in this work is Tensorflow \citep{tensorflow2015-whitepaper}. This library provides lots of tools for machine learning, especially for neural networks. Keras is a wrapper library over Tensorflow and provides easy use of the most common machine learning scenarios \citep{keras}. Tensorflow together with PyTorch \citep{NEURIPS2019_9015} is probably the most frequently used library for deep learning, both providing similar functionality. The reason behind this choice of Tensorflow is the fact, that this thesis builds on the previous work and uses the code developed in Tensorflow.
\subsubsection{Transformers} 
As mentioned in other part of text, this work reuse pretrained language models based on BERT. Transformers library from Hugging Face\citep{Wolf2019HuggingFacesTS} contains many variants of pretrained BERT models and tools for their usage as tokenizers or learning rate schedulers.
\subsubsection{Pandas}
Pandas library \citep{reback2020pandas} serves well for data analysis as it provides data structures like DataFrame, that provides named columns, advanced data indexation, selection, merging, joining, reshaping and other functionality similar to tools provided by e.g.. SQL databases. It does not only provide a rich set of tools but they are also developed with an emphasis on performance optimization.
\subsubsection{Scikit-learn}
Scikit-learn \citep{scikit-learn}  is another useful python library specialized on machine learning. In contrast to tensorflow, scikit-learn focuses on classical machine learning, not on neural networks, providing all important variants of machine learning models as well as supporting tools for training, e.g. cross validation or various metrics.
\subsubsection{Numpy}
Numpy \citep{harris2020array} is an library which provides powerful multidimensional arrays with many predefined operations. It is fast and it is common to use it for numerical operations over number arrays.
\subsubsection{Jupyter Notebook}
Jupyter notebook \citep{jupyter} is a web application for development. In this work, jupyter notebook is used for providing the trained models for exploration. %TODO doplnit odkaz na tu kapitolu
Jupyter suits well for this purpose because it, in addition to a possibility of running a separate parts of code in different cells, also supports visualisations and markdown formatted text and it can be useful especially for explanatory purposes.



\subsection{Code Structure}
Each task has code is in the separate directory as can be seen on picture %TODO udelat obrazek s adresarovou strukturou a barevne oznacit co se ma spoustet
All code for tagging and lemmatization is placed in folder \texttt{morphodita\_research}. Main files are \texttt{morpho\_tagger\_2.py} and \texttt{bert\_finetunning\_simple.py}, which serves for running all experiments relating to tagging and lemmatization. Script arguments are described in more detail in tables \ref{Tab:com_args}, \ref{Tab:mt_com_args} and \ref{Tab:mt2_args}.

%spolecne pro vsechny:
\begin{table}
\centering
\begin{tabular}{ |p{3cm}|p{3,5cm}|p{6cm}| } 
 \hline
 Argument & Values & Description \\ 
 \hline \hline
 accu & int & Accumulation of gradient. Effective batch size is batch\_size * accu.  \\\hline
batch\_size & int & Batch size (without accumulation). \\ \hline
bert & string & Name of the bert model (from huggingFace library) or path to the model.  \\ \hline
  checkp & String & Name of the saved model weights. Saving weights is used instead of the whole model because of some technical issues.  \\ \hline
  debug & 0/1 & Debug mode loads small debug data if available. \\ \hline
  label\_smoothing & decimal number & Coefficient for label smoothing. \\ \hline
  dropout & float &  Dropout amount applied on various places of the network.  \\ \hline
 epochs & "x:l1,y:l2"  & This will perform x epochs with learning rate l1 and y epochs with learning rate l2.   \\ \hline
 layers & None/"att" & If "att", all bert-like model layers are combined with learned weights.  \\ \hline
 warmup\_decay & None /"i:x"/"c:x" & If not None, training will incorporate inverse square root decay or cosine decay for x episodes.  \\ \hline
 fine\_lr & float & Different learning rate for the classification head.  \\ \hline
 \hline
\end{tabular}
\caption{A list of arguments common to all scripts.} 
\label{Tab:com_args}
\end{table}


\begin{table}
\centering
\begin{tabular}{ |p{3cm}|p{3,5cm}|p{6cm}| } 
 \hline
 Argument & Values & Description \\ 
 \hline \hline
 beta\_2 & float & An argument for the optimizer. \\ \hline
 cle\_dim & int & Dimension of character-level embeddings.  \\ \hline
 cont & 0/1 & Evaluating on the test data every epoch.  \\ \hline
 exp & string & Name of logs files.  \\ \hline
 factors & "Lemmas,Tags" & Factors to be predicted -- Lemmas, Tags or both. \\ \hline
word\_dropout & float & Probability of masking a word in the sentence during training.  \\ \hline

\hline

\end{tabular}
\caption{A list of arguments common to both scripts for tagging and lemmatization.} 
\label{Tab:mt_com_args}
\end{table}


\begin{table}
\centering
\begin{tabular}{ |p{3cm}|p{3,5cm}|p{6cm}| } 
 \hline
 Argument & Values & Description \\ 
 \hline \hline
 data & string &  Input data directory. Data are supposed to be divided into train, dev and test \texttt{.txt} files. \\ \hline
 char\_dropout &  float &  Dropout for characters. \\ \hline
embeddings & string & Path to pre-comuputed embeddings for use. \\ \hline
factor\_layers & int & Number of dense-and-dropout blocks for each of factors.  \\ \hline
lemma\_re\_strip &  string & Regular expression for suffix to be stripped from lemma. \\ \hline
lemma\_rule\_min & int & Minimal occurences to keep a lemma.  \\ \hline
predict & string & Script returns only a prediction with model from the path given in this argument.  \\ \hline
rnn\_cell & "LSTM"/"GRU" & Type of rnn cell to use. \\ \hline
rnn\_cell\_dim & int & Dimension for rnn cells.  \\ \hline
rnn\_layers& int & Number of recurrent cell layers.  \\ \hline
we\_dim & int & Dimension of word embeddings.  \\ \hline  


bert\_model & string & Trained checkpoint for loading. Training will continue from this checkpoint. \\ \hline

test\_only & string & Path to the model, which will be load and weights will be printed.  \\ \hline
 \hline
\end{tabular}
\caption{A list of arguments specific to morpho\_tagger\_2.py, with detailed description.} 
\label{Tab:mt2_args}
\end{table}

Sentiment analysis experiments can be run from \texttt{sentiment\_analysis.py} with arguments as described in tables \ref{Tab:com_args} and \ref{Tab:sent_args}.

\begin{table}
\centering
\begin{tabular}{ |p{2cm}|p{4,5cm}|p{6cm}| } 
 \hline
 Argument & Values & Description \\ 
 \hline \hline
 datasets & \{mall,csfd,facebook\} &  Names of the input czech datasets, separated by comma. \\ \hline
 english & float & A percentage of result training data which should be formed by english IMDB dataset.\\ \hline
 
 freeze & {0,1} & 1 means, that bert layers will not be trained. \\ \hline
 seed & int & Inicialization of random seed. \\ \hline
 kfold & "k:i" & Data will be splitted into k folds and i-th fold will be used for test. It serves for running k-fold cross-validation in parallel runs. \\ \hline
 
 \hline
\end{tabular}
\caption{Arguments for \texttt{sentiment\_analysis.py} script.} 
\label{Tab:sent_args}
\end{table}

\subsection{Working example}
The best model is publicly available in the Git repository \footnote{https://github.com/flower-go/DiplomaThesis/tree/master/code/morphodita-research/models\_public} %TODO dodat sentiment model
 and licensed under the Creative Commons Attribution-NonCommercial-ShareAlike 4.0 International Licence \footnote{https://creativecommons.org/licenses/by-nc-sa/4.0/}. If you want to try the model prediction or see a working example of usage of such models, there is public Google Colaboratory \citep{colab} jupyter notebook, which downloads all necessary data and returns predictions for a given text. This notebook si available here: \url{https://colab.research.google.com/github/flower-go/DiplomaThesis/blob/master/PlayWithModels.ipynb}. 








\chapter*{Conclusion}
\addcontentsline{toc}{chapter}{Conclusion}
\label{chap:concl}
In this thesis, we implemented Czech \acrlong{pos} tagging, lemmatization and sentiment analysis with the usage of \acrlong{bert}-like architectures. We achieved \acrlong{sota} results in tagging (accuracy 98.57\%) and lemmatization (accuracy 99.00\%) and joint accuracy of these two tasks 98.19\%, which presents the error reduction of 67\% for tagging accuracy and 53\% for lemmatization accuracy compared to previous public model \citep{Strakova}\footnote{The best publicly available model is to our best knowledge available at \url{https://ufal.mff.cuni.cz/morphodita/users-manual\#czech-morfflex-pdt_model}, which has 95.55\% tagging and 97.96\% lemmatization accuracy.} We also presents new \acrlong{sota} results in sentiment analysis on all three used datasets -- \textit{mall}, \textit{csfd}, \textit{facebook}. We also explored various training techniques and showed the good performance of static embeddings compared to any further training of BERT models. This thesis also examines types of errors BERT helps to solve. All code, text and best models are publicly available on GitHub: \url{https://github.com/flower-go/DiplomaThesis}.

\subsubsection{Future Work}
The BERT-like models ale able to transfer knowledge even across very different languages and also previous work suggests, that training on languages from similar family can improve results in all included languages \citep{Arkhipov2019}, therefore one possible improvement can be in both pre-training the new BERT-like model (similarly to \citep{Straka2021}) on joint dataset for i.e., Czech, Slovak and Polish, or use such multilingual data for task-specific training, for example in sentiment analysis. In addition, results with large XLM-Roberta suggests that selecting this architecture for pretraining could be advantageous despite computational demands. Sentiment analysis data are quite small, although they could be retrieved from pairs (review -- star rating) and many Czech companies have access to such data. Experiments with tagging and lemmatization also shows good results of precomputed embedding approach. The question therefore arises as to whether would not be better to simply add precomputed BERT embeddings into existing sentiment analysis solutions.






%%% Bibliography
%%% Bibliography (literature used as a source)
%%%
%%% We employ bibTeX to construct the bibliography. It processes
%%% citations in the text (e.g., the \cite{...} macro) and looks up
%%% relevant entries in the bibliography.bib file.
%%%
%%% The \bibliographystyle command selects, which style will be used
%%% for references from the text. The argument in curly brackets is
%%% the name of the corresponding style file (*.bst). Both styles
%%% mentioned in this template are included in LaTeX distributions.

%\bibliographystyle{plainnat}    %% Author (year)
\bibliographystyle{mydinat} 
% \bibliographystyle{unsrt}     %% [number]

\renewcommand{\bibname}{Bibliography}

%%% Generate the bibliography. Beware that if you cited no works,
%%% the empty list will be omitted completely.

\bibliography{bibliography, Diplomka}

%%% If case you prefer to write the bibliography manually (without bibTeX),
%%% you can use the following. Please follow the ISO 690 standard and
%%% citation conventions of your field of research.

% \begin{thebibliography}{99}
%
% \bibitem{lamport94}
%   {\sc Lamport,} Leslie.
%   \emph{\LaTeX: A Document Preparation System}.
%   2nd edition.
%   Massachusetts: Addison Wesley, 1994.
%   ISBN 0-201-52983-1.
%
% \end{thebibliography}


%%% Figures used in the thesis (consider if this is needed)
\listoffigures

%%% Tables used in the thesis (consider if this is needed)
%%% In mathematical theses, it could be better to move the list of tables to the beginning of the thesis.
\listoftables

%%% Abbreviations used in the thesis, if any, including their explanation
%%% In mathematical theses, it could be better to move the list of abbreviations to the beginning of the thesis.
\chapwithtoc{List of Abbreviations}

%%% Attachments to the master thesis, if any. Each attachment must be
%%% referred to at least once from the text of the thesis. Attachments
%%% are numbered.
%%%
%%% The printed version should preferably contain attachments, which can be
%%% read (additional tables and charts, supplementary text, examples of
%%% program output, etc.). The electronic version is more suited for attachments
%%% which will likely be used in an electronic form rather than read (program
%%% source code, data files, interactive charts, etc.). Electronic attachments
%%% should be uploaded to SIS and optionally also included in the thesis on a~CD/DVD.
%%% Allowed file formats are specified in provision of the rector no. 72/2017.
\appendix
\chapter{Attachments}

\section{First Attachment}

\openright
\end{document}
